% ==========================================================================
% Section 6 — HRL Architecture: The "Meta-Controller" Approach
% ==========================================================================
\section{Architecture 2: Hierarchical RL with an LLM Meta-Controller}
\label{ch:hrl_architecture}

The previous chapter established the \emph{Mixing Board} architecture, representing a "soft integration" of Large Language Models into MARL. In that paradigm, the LLM acted as an environmental architect, manipulating the gravitational pull of potential-based rewards to gently guide the agents toward logical milestones. While this separation preserves the underlying Markov properties, it restricts the LLM to an indirect role: it can only suggest a direction, not enforce a strategy. 

To fully explore the spectrum of neuro-symbolic integration, we must ask: what happens when we grant the LLM direct semantic authority over the agents' cognitive state?

This section details the second integration strategy: a Hierarchical Reinforcement Learning (HRL) framework. Here, we transition from soft environmental guidance to hard architectural integration. The LLM ceases to be a distant architect and becomes a direct \emph{Meta-Controller}. Instead of adjusting continuous weights over a fixed greedy planner, the Meta-Controller replaces the planner entirely. It selects discrete high-level \emph{Options} (macro-actions) that strictly govern the behavior of the low-level MAPPO policy.

In our established metaphor, the Commander is no longer just shifting the gravity; it is issuing explicit, discrete commands directly to the sensory inputs of the lumberjack and miner.

The exposition follows the implementation order:
\begin{enumerate}
    \item \textbf{Option Definitions} (\cref{sec:option_defs}): Defining the discrete macro-actions available to the Commander.
    \item \textbf{Discrete Option Selection} (\cref{sec:discrete_selection}): How the LLM is triggered to assign these explicit tasks.
    \item \textbf{One-Hot Observation Augmentation} (\cref{sec:onehot_aug}): How the Commander's text-based orders are mathematically injected into the MAPPO neural network.
    \item \textbf{Bipartite Reward Structure} (\cref{sec:bipartite_reward}): The dual-reward system required to train the workers to follow these new explicit orders.
    \item \textbf{Staleness and Counterfactual Reflection} (\cref{sec:lifecycle_guardrails}): How the Commander diagnoses its own mistakes when an order fails on the ground.
\end{enumerate}

% ──────────────────────────────────────────────────────────────────────────
\subsection{System Overview}
\label{sec:hrl_overview}

\Cref{fig:hrl_system_flow} illustrates this two-level command architecture. The Meta-Controller (LLM) observes inventory transitions and assigns explicit Options to each agent asynchronously. Unlike the previous approach where the MAPPO policy remained blind to the LLM's existence, here, the low-level MAPPO policy receives the active option directly as part of its mathematical observation space. This creates a true, tightly coupled two-level decision hierarchy.

\begin{figure}[H]
\centering
\begin{tikzpicture}[
    box/.style={draw, rounded corners, minimum width=2.6cm, minimum height=0.9cm,
                font=\small\sffamily, align=center, fill=gray!8},
    llmbox/.style={draw, rounded corners, minimum width=2.6cm, minimum height=0.9cm,
                font=\small\sffamily, align=center, fill=blue!8},
    arr/.style={-{Stealth[length=2.5mm]}, thick},
    node distance=1.6cm and 2cm
]
    % Top row: LLM level
    \node[llmbox] (llm) {LLM\\Meta-Controller};
    \node[box, right=3cm of llm] (ctrl) {OptionController\\$o_t^{(0)}, o_t^{(1)}$};

    % Bottom row: MAPPO level
    \node[box, below=2cm of llm] (env) {Env Step\\$s_t \to s_{t+1}$};
    \node[box, right=3cm of env] (agent) {MAPPO Agent\\$\pi_\theta(a \mid s, o)$};

    % LLM -> Controller
    \draw[arr] (llm) -- node[above,font=\footnotesize]{JSON} (ctrl);

    % Controller -> Agent (one-hot)
    \draw[arr] (ctrl.south) -- ++(0,-0.6) -|
        node[near start,right,font=\footnotesize]{one-hot $\mathbf{e}_{o}$} (agent.north);

    % Agent -> Env
    \draw[arr] (agent.south) -- ++(0,-0.4) -|
        node[near start,below,font=\footnotesize]{$a_t$} (env.south);

    % Env -> Agent (reward)
    \draw[arr] (env.east) -- node[below,font=\footnotesize]{$r_t^{\text{env}} + r_t^{\text{intr}}$} (agent.west);

    % Env -> LLM (trigger)
    \draw[arr,dashed] (env.north) -- node[left,font=\footnotesize]{\scriptsize success / stale} (llm.south);

    % Env -> Controller (reset on terminal)
    \draw[arr,dashed] (env.north east) -- node[right,font=\footnotesize,pos=0.3]{\scriptsize terminal} (ctrl.south west);
\end{tikzpicture}
\caption[Two-Level HRL Architecture]{Two-level HRL architecture. The LLM Meta-Controller is triggered by option success or staleness (dashed). The OptionController passes one-hot embeddings to the MAPPO agent at every step (solid), directly altering the agent's cognitive state.}
\label{fig:hrl_system_flow}
\end{figure}


% ──────────────────────────────────────────────────────────────────────────
\subsection{Option Definitions}
\label{sec:option_defs}

In the dynamic shaping approach, the LLM merely influenced a hard-coded planner. In this Hierarchical Reinforcement Learning (HRL) architecture, the LLM completely replaces the planner. It selects from a finite set of explicit macro-actions, or \emph{Options}, and directly commands the low-level MAPPO agents.

An Option $o\in\mathcal{O}$ is mathematically defined by the triple $\langle \mathcal{I}_o, \pi_o, \beta_o \rangle$: an initiation set, an intra-option policy, and a termination condition~\parencite{Sutton_Precup_Singh_1999}.
In our framework, the option set $\mathcal{O}$ is derived directly from the dependency DAG, providing 10 discrete macro-actions as detailed in Table~\ref{tab:option_props}.

\begin{table}[H]
\centering
\caption{Option properties: target zone, responsible agent, and initiation condition.}
\label{tab:option_props}
\small
\begin{tabular}{@{}lccc@{}}
\toprule
\textbf{Option} & \textbf{Zone} & \textbf{Agent} & \textbf{Initiation Condition} \\
\midrule
\texttt{COLLECT\_WOOD}  & Wood (0)      & 0     & Always \\
\texttt{COLLECT\_STONE} & Stone (1)     & 1     & Always \\
\texttt{CRAFT\_PICKAXE} & Workbench (2) & Both  & wood $\ge 1$, stone $\ge 1$, pickaxe $< 1$ \\
\texttt{MINE\_IRON}     & Iron (3)      & 1     & pickaxe $\ge 1$ \\
\texttt{CRAFT\_SWORD}   & Workbench (2) & Both  & iron $\ge 1$, sword $< 1$ \\
\texttt{CRAFT\_ARMOR}   & Workbench (2) & Both  & iron $\ge 1$, armor $< 1$ \\
\texttt{BUILD\_BRIDGE}  & Bridge (4)    & Both  & wood $\ge 1$, bridge $< 1$ \\
\texttt{FIGHT\_ENEMY}   & Enemy (5)     & Both  & bridge, sword, armor $\ge 1$ \\
\texttt{COLLECT\_GOLD}  & Gold (6)      & 1     & enemy defeated, gold $< 1$ \\
\texttt{IDLE}           & ---           & Both  & Always \\
\bottomrule
\end{tabular}
\end{table}

Rather than building 10 separate neural networks, the intra-option policy $\pi_o$ is executed by the single, shared MAPPO network, which conditions its kinetic movements based on the active option (detailed in Section~\ref{sec:onehot_aug}).

The termination condition ($\beta_o$) is not based on spatial coordinates, but on semantic success. For example, if the Commander radios the \texttt{COLLECT\_WOOD} Option to the lumberjack, that Option only terminates when the system registers an increase in the wood inventory.

% ──────────────────────────────────────────────────────────────────────────
\subsection{Discrete Option Selection by the LLM}
\label{sec:discrete_selection}

Because the Meta-Controller issues explicit, discrete commands, it does not need to continuously monitor the mathematical variance of the Critic network (unlike the Mixing Board approach). Instead, the Commander acts reactively, stepping in only when a worker completes a task or fails to make progress.

\subsubsection{Trigger Conditions: Success and Staleness}

Option selection is triggered dynamically under two specific conditions, evaluated per-en\-vi\-ron\-ment:

\begin{enumerate}
    \item \textbf{Success:} The agent successfully acquired the item (the inventory slot changed), meaning the current Option has terminated and the agent awaits a new order.
    \item \textbf{Staleness:} The active option has been running for more than $T_{\text{stale}}=150$ steps without completion.
    In MARL, a step $t \to t{+}1$ represents a physical movement or action in the grid. If an agent executes 150 low-level kinetic actions without satisfying its macro-objective, the system assumes the agent is stuck (e.g., trying to mine iron without a pickaxe) and forces an LLM intervention.
\end{enumerate}

Environments that are already awaiting an LLM response (\texttt{llm\_pending}) or are on a brief 50-step cooldown are safely excluded from triggering new queries.

\subsubsection{Prompt Structure}

When a trigger fires, the system translates the grid-world's physical inventory state into a text-based JSON schema.
For a standard \emph{Success} trigger, the Commander receives a Chain-of-Thought prompt~\parencite{Wei_Wang_Schuurmans_Bosma_Ichter_Xia_Chi_Le_Zhou_2022} detailing the current inventory and the DAG dependencies, requesting exactly one discrete Option per agent:

\begin{hegquote}
\begin{verbatim}
You are the high-level Cognitive Orchestrator for two MARL agents.
Strict Dependency DAG: (Wood+Stone) -> Pickaxe -> Iron
    -> (Sword+Armor) -> Bridge -> Enemy -> Gold.

### CURRENT STATE:
Inventory: {inventory}
Agent 0 Status: {a0_status}
Agent 1 Status: {a1_status}

### YOUR TASK:
Assign exactly ONE discrete Option to each agent.
Available Options: ["COLLECT_WOOD", "COLLECT_STONE", "CRAFT_PICKAXE",
    "MINE_IRON", "CRAFT_SWORD", "CRAFT_ARMOR", "BUILD_BRIDGE",
    "FIGHT_ENEMY", "COLLECT_GOLD", "IDLE"]

Respond ONLY with valid JSON:
{
  "dag_check": "<1 sentence reasoning>",
  "agent_0_option": "<Option>",
  "agent_1_option": "<Option>"
}
\end{verbatim}
\end{hegquote}

\subsubsection{Asynchronous Query Management}

Training MARL across 128 parallel environments generates significant query volume. To optimize throughput, the system deduplicates queries by grouping environments with identical inventory states, sending a single prompt per unique state to the LLM, and broadcasting the response to all matching environments.

Additionally, because LLM inference takes several seconds, continuous environment execution would cause \textbf{state drift} (latency-induced non-stationarity)~\parencite{Katsikopoulos_Engelbrecht_2003}, rendering the LLM's eventual decision outdated. Following established LLM-agent practices~\parencite{Zhu_etal_2023}, any environment awaiting a response is paused. This ensures the Meta-Controller's decision is applied to the exact state it observed, maintaining strict temporal alignment between strategy and execution.

% ──────────────────────────────────────────────────────────────────────────
\subsection{One-Hot Observation Augmentation}
\label{sec:onehot_aug}

Once the Commander has selected a discrete Option, that command must be physically transmitted to the worker. In the MAPPO framework, this is achieved by appending a one-hot vector directly into the agent's localized observation space, effectively acting as a ``radio receiver'' for the Commander's orders. We can see how exactly in Equation \ref{eq:hrl_obs}.

\subsubsection{Dynamic Enemy State}

Unlike the Mixing Board approach, which utilized the baseline static environment, the HRL architecture operates in an extended environment that introduces a \emph{dynamic enemy}. In this variant, the enemy is no longer a static zone interaction. Instead, the enemy tracks the agents' positions and actively chases any agent that crosses the river, creating a real-time pursuit challenge. The combat system is designed to reward coordination: if both agents attack the enemy simultaneously, they deal 50~HP of damage per step, compared to only 10~HP for a solo attack. While this 5$\times$ coordination multiplier was intended to make the \texttt{FIGHT\_ENEMY} option inherently cooperative, the overarching time penalty ($-0.01$ per step) often creates conflicting incentives. As explored in Chapter~\ref{ch:experiments}, the agents occasionally discover that having one agent rush the enemy while the other crafts is more time-efficient than waiting to launch a coordinated assault.

To give the agents spatial awareness of this dynamic threat, the three dimensions that were occupied by the goal embedding in the Mixing Board approach (indices 2--4) are repurposed to carry the enemy's real-time state: its $(x, y)$ position and its normalized health $h/100$. Providing this state information is standard practice in fully observable Markov Decision Processes (MDPs); it simply represents the agents' physical line-of-sight, which is required for the low-level PPO policy to learn the kinetic micro-actions of pursuit and combat.

\subsubsection{Observation Construction}

The vector input to the MAPPO network in this HRL mode is expanded to a 25-dimensional vector:
\begin{equation}
\label{eq:hrl_obs}
    \mathbf{v}_t^{(i)}
    =
    \left[
        \underbrace{\mathbf{p}^{(i)}}_{\mathbb{R}^2}
        \;\Big\|\;
        \underbrace{\mathbf{d}^{\text{enemy}}}_{\mathbb{R}^3}
        \;\Big\|\;
        \underbrace{\mathbf{x}}_{\mathbb{R}^{10}\text{ (inv)}}
        \;\Big\|\;
        \underbrace{\mathbf{e}_{o_t^{(i)}}}_{\mathbb{R}^{10}\text{ (one-hot)}}
    \right]
    \;\in\; \mathbb{R}^{25}
\end{equation}
where $\mathbf{p}^{(i)}$ is the agent's spatial position, $\mathbf{d}^{\text{enemy}} = (x_{\text{enemy}},\, y_{\text{enemy}},\, h_{\text{enemy}}/100)$ is the dynamic enemy state vector, $\mathbf{x}$ is the shared inventory vector, and $\mathbf{e}_{o}\in\{0,1\}^{10}$ is the one-hot encoding of the active option.

For example, if an agent is located at coordinates $(3, 4)$, the enemy is at $(5, 5)$ with 50\% health, the inventory contains exactly 1 wood, and the Commander has assigned the first option (e.g., \texttt{COLLECT\_WOOD}, index 0), the augmented vector is constructed as follows:
\begin{equation}
    \mathbf{v}_t^{(i)}
    =
    \left[
        \underbrace{3,\; 4}_{\mathbf{p}^{(i)}}
        \;\Big\|\;
        \underbrace{5,\; 5,\; 0.5}_{\mathbf{d}^{\text{enemy}}}
        \;\Big\|\;
        \underbrace{1,\; 0,\dots, 0}_{\mathbf{x} \in \mathbb{R}^{10}}
        \;\Big\|\;
        \underbrace{1,\; 0,\dots, 0}_{\mathbf{e}_o \in \mathbb{R}^{10}}
    \right]
\end{equation}

This design expands the $\mathbb{R}^{15}$ observation space of the Mixing Board approach to $\mathbb{R}^{25}$. Rather than increasing the spatial observation dimensionality, the three dimensions previously occupied by the continuous goal embedding are efficiently repurposed to carry the enemy state. The additional 10 dimensions are strictly dedicated to the one-hot option vector.

By appending this one-hot vector $\mathbf{e}_{o_t^{(i)}}$ directly to the observation space, the MAPPO network learns an \emph{option-conditional policy} $\pi_\theta(a \mid s, o)$. This one-hot encoding acts as a mathematical switch, allowing a single neural network to dynamically shift its behavior based on the Commander's active assignment, avoiding the need to train a separate policy network for each of the 10 distinct options.

% ──────────────────────────────────────────────────────────────────────────
\subsection{Bipartite Reward Structure}
\label{sec:bipartite_reward}

In standard RL, the environment reward is the ground truth. However, in our sparse environment, relying solely on this is analogous to receiving feedback only at the very end of a long episode. To ensure the agents learn the day-to-day behaviors required to achieve that final milestone, the HRL approach uses a bipartite composite reward: a sparse environment reward and a dense intrinsic reward.

\subsubsection{Extrinsic and Intrinsic Components}
\label{sec:intrinsic_reward}

As established in Section~\ref{sec:env_design}, the extrinsic environment reward $r_t^{\text{env}}$ is the true, unmodified sparse signal from the grid-world (refer to Table~\ref{tab:env_rewards}). It consists of a constant time-step penalty of $-0.01$ to encourage speed, plus massive milestone bonuses upon successfully completing a physical subtask.

Conversely, the intrinsic (shaping) reward is a dense synthetic shaping signal generated internally by the architecture to keep the agent motivated toward the Commander's assigned target. 
It consists of two sub-components: a continuous distance-based shaping mechanism, and a discrete option success bonus.

First, the continuous distance-based shaping reward toward the active option's target zone is computed identically to the PBRS framework~\parencite{Ng_Harada_Russell_1999} established in Section~\ref{subsec:pbrs}, but applied to the option's target rather than the planner's target:
\begin{equation}
\label{eq:hrl_shaping}
    r_t^{\text{shape},(i)}
    \;=\;
    \mathds{1}_{\{o^{(i)} \neq \texttt{IDLE}\}}
    \;\cdot\;
    \lambda_{\text{intr}}
    \bigl(\gamma\,\Phi^{(i)}(s_{t+1}, z_o) - \Phi^{(i)}(s_t, z_o)\bigr)
\end{equation}
where $z_o$ is the target zone of option~$o$ (\cref{tab:option_props}), $\Phi$ is the potential function, and $\lambda_{\text{intr}} = 0.05$ is the shaping scaling factor (equivalent to $\lambda_{\text{scale}}$ from the Mixing Board approach).

As established in previous chapters, this canonical difference form $\gamma\,\Phi(s_{t+1}) - \Phi(s_t)$ guarantees policy invariance and naturally prevents the agents from exploiting the reward signal by stepping back and forth~\parencite{Ng_Harada_Russell_1999}.

The indicator function $\mathds{1}_{\{o^{(i)} \neq \texttt{IDLE}\}}$ acts as a coordination mask. An agent assigned the \texttt{IDLE} option is actively waiting for its partner to satisfy a dependency; applying distance-based shaping here would incorrectly penalize it for remaining stationary. The mask zeroes out this shaping signal, allowing the agent to safely wait without destabilizing the value network.

Second, when an agent completes its active option (the relevant inventory slot increases), the distance shaping is overridden, and it receives a flat intrinsic success bonus of $+1.0$:
\begin{equation}
\label{eq:success_bonus}
    r_t^{\text{intr},(i)}
    =
    \begin{cases}
        1.0 & \text{if option } o^{(i)} \text{ terminated successfully} \\[2pt]
        r_t^{\text{shape},(i)} \text{ (\cref{eq:hrl_shaping})} & \text{otherwise}
    \end{cases}
\end{equation}

Success is detected by comparing the inventory before and after the environment step.

\subsubsection{Composite Reward}

The total reward signal fed to the PPO algorithm is a weighted sum of the environment diploma and the intrinsic breadcrumbs:
\begin{equation}
\label{eq:hrl_total_reward}
\boxed{
    r_t^{\text{total}}
    = r_t^{\text{env}}
    + \alpha_{\text{intr}} \cdot r_t^{\text{intr}}
}
\end{equation}
with $\alpha_{\text{intr}} = 0.5$. This coefficient was chosen to balance the dense intrinsic signal against the sparse environmental milestones; too large a value causes the policy to over-optimise for intrinsic shaping at the expense of actual task completion.

Note the key difference from the Mixing Board approach: here, no $\lambda$ decay schedule is applied. The intrinsic coefficient remains fixed because the LLM continuously updates the \emph{direction} (which option to pursue) rather than the magnitude of shaping.

% ──────────────────────────────────────────────────────────────────────────
\subsection{Mathematical Impact on the MAPPO Objective}
\label{sec:hrl_math_impact}

Unlike the Dynamic approach, which acted purely as an environment architect through the reward signal, the Meta-Controller strikes the MAPPO equations in \emph{two places simultaneously}: the observation space (Actor) and the reward space (Critic).

On the Actor side, the one-hot option conditioning transforms the decentralized policy of agent $i$ from $\pi_\theta(a_t^{(i)} \vert{} s_t^{(i)})$ into an option-conditional distribution $\pi_\theta(a_t^{(i)} \vert{} s_t^{(i)}, o_t^{(i)})$. 
The Actor's goal in PPO is to maximize the probability of good actions, but it uses a ``clipped objective'' to ensure the network does not update its weights too aggressively in a single step. 
By injecting the discrete option $o_t^{(i)}$, the probability ratio $r_t^{(i)}(\theta)$ between the new and old policy becomes option-aware:
\begin{equation}
    r_t^{(i)}(\theta) = \frac{\pi_\theta(a_t^{(i)} \vert{} s_t^{(i)}, o_t^{(i)})}{\pi_{\theta_{\text{old}}}(a_t^{(i)} \vert{} s_t^{(i)}, o_t^{(i)})}
\end{equation}
and the clipped surrogate objective updates accordingly:
\begin{equation}
    L^{\text{CLIP}}(\theta) = \mathbb{E} \left[ \min\left( r_t^{(i)}(\theta) A_t^{(i)},\; \text{clip}\bigl(r_t^{(i)}(\theta), 1 - \epsilon, 1 + \epsilon\bigr) A_t^{(i)} \right) \right]
\end{equation}

On the Critic side, learning is driven by the Advantage function ($A_t^{(i)}$). Remember that the Advantage measures how much better a chosen action was compared to the baseline expectation. This baseline is provided by the Centralized Value function $V_\phi$. 
To maintain the CTDE paradigm established in Section~\ref{subsec:ctde}, the Critic evaluates the \emph{global} environment state $\mathbf{s}_t$. In this HRL architecture, it also receives the joint options of all agents $\mathbf{o}_t = (o_t^{(0)}, \dots, o_t^{(N)})$. The bipartite reward feeds into this option-conditional Advantage:
\begin{equation}
    A_t^{(i)} = r_t^{\text{total}, (i)} + \gamma V_\phi(\mathbf{s}_{t+1}, \mathbf{o}_{t+1}) - V_\phi(\mathbf{s}_t, \mathbf{o}_t)
\end{equation}
Because the Value function $V_\phi$ now receives the globally augmented state $\mathbf{s}_t' = [\mathbf{s}_t \parallel \mathbf{e}_{\mathbf{o}_t}]$, the Critic learns to predict entirely different baseline expected returns depending on the combination of options actively guiding the agents.

The consequence is a tightly coupled learning loop: the Actor gradient explicitly links the local physical state $s_t^{(i)}$, the Commander's specific order $o_t^{(i)}$, and the intrinsic success of that command through $A_t^{(i)}$. The network learns distinct sub-policies for each option, not through separate networks, but through a single policy that conditions its behavior on which option the LLM has assigned.

% ──────────────────────────────────────────────────────────────────────────
\subsection{Execution Lifecycle and Stability Guardrails}
\label{sec:lifecycle_guardrails}

To ensure the HRL architecture remains robust over millions of training steps, several lifecycle management and stability mechanisms are integrated directly into the environment loop.

\subsubsection{Episode Initialization and Reset}
When an environment reaches a terminal state (gold mined, game over, or truncation at 800 steps), the environment is hard-reset. The agents are assigned initial options corresponding to the first layer of the DAG (\texttt{COLLECT\_WOOD} for the lumberjack, \texttt{COLLECT\_STONE} for the miner), ensuring they begin collecting resources immediately upon episode start without waiting for an initial LLM query.

\subsubsection{Staleness Detection and Counterfactual Reflection}
During execution, an internal watchdog tracks the age (number of steps) of each active option. If an option exceeds $T_{\text{stale}}=150$ steps without termination, the environment is marked as stale and the LLM is re-queried. The \texttt{IDLE} options are excluded from this trigger, as intentional waiting is a valid cooperative strategy.

When a staleness trigger fires, the LLM receives a \emph{counterfactual reflection prompt} instead of the standard assignment prompt. Rather than simply asking for a new assignment, this prompt forces the LLM to diagnose \textbf{why} its previous assignment failed (e.g., a missing DAG precondition like \texttt{MINE\_IRON} without a pickaxe) before issuing new options:

\begin{hegquote}
\begin{verbatim}
### REFLECTION REQUIRED:
Agent 0 was assigned "{a0_option}" and Agent 1 was assigned "{a1_option}".
After {stale_steps} steps, NEITHER has succeeded.

### YOUR TASK:
1. First, diagnose what prerequisite is MISSING.
2. Then assign corrected options that respect the DAG.

Respond ONLY with JSON:
{
  "reflection": "<diagnosis>",
  "dag_check": "<verification>",
  "agent_0_option": "<Option>",
  "agent_1_option": "<Option>"
}
\end{verbatim}
\end{hegquote}

This mechanism leverages the LLM's zero-shot reasoning to debug its own prior hallucinations, synergizing reasoning with counterfactual critique~\parencite{Yao_Zhao_Yu_Du_Shafran_Narasimhan_Cao_2023, Shinn_Cassano_Berman_Gopinath_Narasimhan_Yao_2023}.

\subsubsection{Target KL Divergence Guard}
To protect the policy network from destructive updates when intrinsic reward signals shift abruptly after an LLM reassignment, the training loop employs a target KL divergence early-stopping mechanism~\parencite{Schulman_Wolski_Dhariwal_Radford_Klimov_2017}. During each PPO epoch, the approximate KL divergence between the old and new policies is monitored:
\begin{equation}
\label{eq:approx_kl}
    D_{\text{KL}}^{\text{approx}}
    = \mathbb{E}\bigl[-\log \pi_{\theta_{\text{new}}}(a \mid s, o)
      + \log \pi_{\theta_{\text{old}}}(a \mid s, o)\bigr]
\end{equation}
If $D_{\text{KL}}^{\text{approx}} > 0.015$ (an empirical threshold widely adopted in standard PPO implementations), the remaining PPO epochs are aborted for that update, ensuring stable policy convergence.

% ──────────────────────────────────────────────────────────────────────────
\subsection{Comparison with the Mixing Board Approach}
\label{sec:hrl_vs_dyn}

Before detailing the specific fine-tuning pipeline required for the Meta-Controller, it is instructive to contrast this architecture with the Dynamic Reward Shaping (Mixing Board) approach developed in Chapter 5. While both paradigms seek to integrate LLM reasoning into MARL, they operate at fundamentally different levels of abstraction. 

The Mixing Board is a \emph{soft integration} strategy: the LLM gently nudges the priority of intrinsic rewards, but the MARL agents remain entirely responsible for discovering the physical behaviors. Conversely, the Meta-Controller represents a \emph{hard integration}: the LLM acts as a strict supervisor, dictating explicit macro-actions (options) that the agents must execute. This shift from continuous priority logits to discrete option assignments necessitates significant structural changes, particularly regarding LLM query frequency, observation augmentation, and safety mechanisms. These core architectural differences are summarized in Table~\ref{tab:approach_comparison}.

\begin{table}[H]
\centering
\caption{Structural comparison of the two LLM integration strategies.}
\label{tab:approach_comparison}
\small
\begin{tabular}{@{}p{3.5cm}p{5cm}p{5cm}@{}}
\toprule
\textbf{Aspect} & \textbf{Mixing Board (Soft Integration)} & \textbf{Meta-Controller (Hard Integration)} \\
\midrule
LLM output space &
    Continuous logits $\in \mathbb{R}^7$ &
    Discrete options $\in \mathcal{O}^2$ \\
LLM controls &
    Subtask priority weights &
    Explicit per-agent macro-actions \\
Trigger mechanism &
    Plateau + Z-score on TD error &
    Option success or staleness \\
Query granularity &
    Global (once per trigger) &
    Per-environment (batched) \\
Observation augmentation &
    None (blind to LLM) &
    One-hot option $\in \{0,1\}^{10}$ \\
Reward structure &
    PBRS with $\lambda$ decay &
    Bipartite: sparse env + dense intrinsic (fixed $\alpha$) \\
LLM query frequency &
    Low (plateau-triggered) &
    High (per-option completion) \\
Safety mechanism &
    Softmax budget, EMA, rollback &
    DAG in prompt, reflection, env freezing \\
\bottomrule
\end{tabular}
\end{table}

% ──────────────────────────────────────────────────────────────────────────
\subsection{QLoRA Fine-Tuning Pipeline}
\label{sec:qlora_pipeline}

While Section~\ref{sec:qlora} established the theoretical foundations of Quantized Low-Rank Adaptation, this section details the complete practical pipeline used to produce the specialized Commander adapter deployed throughout the HRL experiments evaluated in Chapter~\ref{ch:experiments}. As will be demonstrated, a surprisingly brief fine-tuning process on a modest synthetic dataset is sufficient to dramatically improve the Meta-Controller's decision quality.

\subsubsection{Oracle Dataset Generation}

The training dataset is generated synthetically to ensure exhaustive coverage of the crafting DAG. Exploratory successful runs from the agents helped build the dataset. A deterministic oracle systematically enumerates all meaningful inventory combinations and applies hardcoded logic to determine the optimal \texttt{(agent\_0\_option, agent\_1\_option)} assignment for each state. For example: ``if no pickaxe and wood $\ge 1$ and stone $\ge 1$, then both agents should \texttt{CRAFT\_PICKAXE}.''

To prevent the model from over-indexing on trivially simple early-game states, the generated states are grouped into six progression phases (e.g., early collection, iron mining, endgame combat). Each phase is upsampled equally to ensure a balanced training signal.

Two augmentation techniques are applied to improve generalization:

\begin{enumerate}
    \item \textbf{Stochastic Oracle Augmentation (SOA):} To prevent the model from overfitting to the perfectly logical inventory progressions generated by the script, a randomly selected inventory field is zeroed out with a 10\% probability before querying the oracle. For example, a valid mid-game state like \texttt{[iron: 1, pickaxe: 1]} might be corrupted to the physically impossible state \texttt{[iron: 1, pickaxe: 0]}. By passing this corrupted state to the oracle to generate the training label, the LLM is explicitly trained on how to gracefully recover from weird, out-of-distribution inventory combinations. This robustness is critical because observation delays and asynchronous race conditions during live RL rollouts frequently produce these mathematically inconsistent states.
    \item \textbf{Reflection Augmentation (40\% injection rate):} A synthetic hallucinated failure is prepended to the prompt: ``\emph{You previously chose [Incorrect Option]. It failed after [N] steps.}'' The model must then output the correct option despite this misleading prior. While intended to teach the Commander to recover from logical errors, Section~\ref{sec:kl_guard} demonstrates that this training signal inadvertently caused the model to assume \emph{any} reported failure implies a logical mistake (the ``Reflection Contradiction'').
\end{enumerate}

The final raw dataset contains \textbf{1,200 balanced examples} spanning all DAG phases and agent status combinations.

\subsubsection{ChatML Formatting and Dataset Splits}

The raw prompt and completion pairs are converted into Qwen's native ChatML format~\parencite{Yang_Li_Yang_Zhang_Hui_Zheng_Yu_Gao_Huang_Lv_et_al._2025}, structuring each example as a three-turn conversation (\texttt{system}, \texttt{user}, \texttt{assistant}). The system message contains the Commander's role description and DAG rules, the user message provides the current inventory state (and, for reflection examples, the failure context), and the assistant message contains the expected JSON output.

After a 90/10 train--validation split and a $2\times$ training set duplication (to ensure sufficient gradient updates given the small dataset), the final corpus comprises \textbf{2,160 training examples} and \textbf{120 validation examples}.

\subsubsection{Training Configuration}

Fine-tuning was performed on the local workstation's NVIDIA RTX 5070 GPU (12GB VRAM), as described in Section~\ref{sec:env_hardware}. The strict memory envelope forced several architectural compromises documented in Table~\ref{tab:qlora_training_config}.

\begin{table}[H]
    \centering
    \caption[QLoRA Fine-Tuning Hyperparameters]{QLoRA fine-tuning hyperparameters. Values marked with $\dagger$ were constrained by the 12GB VRAM budget.}
    \label{tab:qlora_training_config}
    \begin{tabular}{@{}llr@{}}
        \toprule
        \textbf{Category} & \textbf{Parameter} & \textbf{Value} \\
        \midrule
        Base Model & Architecture & Qwen 2.5 7B-Instruct \\
        & Source & HuggingFace Hub~\parencite{Wolf_Debut_Sanh_Chaumond_Delangue_Moi_Cistac_Rault_Louf_Funtowicz_et_al._2020} \\
        \midrule
        Quantization & Precision & 4-bit NormalFloat (NF4) \\
        & Compute dtype & bfloat16 \\
        & Double quantization & Enabled \\
        \midrule
        LoRA Adapter & Rank $r$ & 16 \\
        & Scaling $\alpha$ & 32 \\
        & Dropout & 0.05 \\
        & Target modules & All 7 (q/k/v/o/gate/up/down) \\
        \midrule
        Training & Epochs & 1 \\
        & Batch size$^\dagger$ & 1 \\
        & Gradient accumulation & 8 (eff.\ batch = 8) \\
        & Learning rate & $2 \times 10^{-4}$ \\
        & Scheduler & Cosine with warmup (10\%) \\
        & Optimizer & Paged AdamW 8-bit~\parencite{bitsandbytes} \\
        & Max sequence length$^\dagger$ & 1,024 tokens \\
        & Gradient checkpointing$^\dagger$ & Enabled \\
        \midrule
        Infrastructure & Library & HuggingFace TRL~\parencite{trl} \\
        & Adapter library & PEFT v0.19.1~\parencite{peft} \\
        & Logging & TensorBoard \\
        \bottomrule
    \end{tabular}
\end{table}

A critical training detail is the application of \emph{response-only loss masking}, implemented via the TRL library~\parencite{trl}. By default, causal language models learn by predicting every subsequent word in a sequence, meaning they would normally compute an error gradient for the prompt itself. This wastes training capacity on memorizing instructions. To solve this, a custom data collator leverages PyTorch's native \texttt{ignore\_index} mechanic: it sets the target labels of all system and user prompt tokens to $-100$. Mathematically, the Cross-Entropy loss function ignores any token with a $-100$ label, computing zero gradient for it. For example, in the sequence \texttt{[User] State is X [Assistant] \{"opt": 1\}}, every token preceding the JSON response is assigned a $-100$ label, while only the JSON characters receive valid target IDs. As a result, the model is penalized exclusively on its final JSON output, ensuring that every weight update directly improves its strategic decision-making rather than memorizing boilerplate text.

\subsubsection{Training Results and Adapter Deployment}

The complete fine-tuning process converged in approximately \textbf{3.1 hours} (1,910 steps) on the RTX 5070. As shown in Figure~\ref{fig:qlora_training}, the training loss dropped rapidly from 1.605 to 0.165, with the final evaluation metrics confirming near-perfect task acquisition:

\begin{itemize}
    \item Evaluation loss: \textbf{0.0007}
    \item Mean token accuracy: \textbf{99.98\%}
\end{itemize}

\begin{figure}[H]
    \centering
    \includegraphics[width=0.85\textwidth]{../../plots/qlora/QLoRA_Training_Performance.png}
    \caption[QLoRA Adapter Fine-Tuning Performance]{QLoRA adapter fine-tuning performance over 1,910 training steps. The training loss (left axis) converges rapidly within the first 500 steps, while the mean token accuracy (right axis) saturates near 100\%.}
    \label{fig:qlora_training}
\end{figure}

The resulting LoRA adapter weighs only \textbf{154~MB}, approximately 1\% of the full 15~GB base model, yet encodes the complete DAG reasoning and JSON formatting capabilities required by the Meta-Controller. For deployment during RL training, the adapter is loaded via the PEFT library and merged with the base model in memory, or alternatively converted to GGUF format and served through the Ollama inference framework (Section~\ref{sec:env_hardware}). Both pathways produce inference latencies compatible with the synchronous RL training loop, as quantified in Section~\ref{sec:computational_overhead}. This fine-tuned adapter serves as the Meta-Controller for all ``QLoRA'' HRL configurations evaluated in Chapter~\ref{ch:experiments}.

% ──────────────────────────────────────────────────────────────────────────
\subsection{Summary}
\label{sec:hrl_summary}

The Meta-Controller approach establishes a strict separation of concerns: the LLM Commander dictates \emph{what to do} (strategic option selection), while the MAPPO policy learns \emph{how to do it} (kinetic execution). By appending a one-hot option encoding to the state vector, a single policy network learns option-specific behaviors. Simultaneously, the bipartite reward structure ensures the policy receives both dense intrinsic guidance and sparse environmental grounding.

Crucially, this architecture functions as an automated curriculum generator. Unlike traditional Curriculum Learning (Section~\ref{sec:curriculum_learning}), which requires manually designing environments of escalating difficulty, the Meta-Controller orchestrates a skill curriculum dynamically. It guides agents from basic resource gathering to advanced combat coordination entirely within a single, static environment.

\begin{table}[H]
\centering
\caption{Key hyperparameters of the Meta-Controller approach.}
\label{tab:hrl_hyperparams}
\begin{tabular}{@{}lll@{}}
\toprule
\textbf{Component} & \textbf{Parameter} & \textbf{Value} \\
\midrule
Options          & $|\mathcal{O}|$                & 10 \\
Staleness        & $T_{\text{stale}}$             & 150 steps \\
Cooldown         & Post-query cooldown            & 50 steps \\
Intrinsic reward & Distance scale $\lambda_{\text{intr}}$ & 0.05 \\
                 & Success bonus                  & $+1.0$ \\
Composite reward & $\alpha_{\text{intr}}$         & 0.5 \\
                 & $D_{\max}$                     & 85.0 \\
Target KL        & KL threshold                   & 0.015 \\
Network          & One-hot dimension              & 10 \\
                 & Total \texttt{vec\_dim}        & 25 \\
\bottomrule
\end{tabular}
\end{table}
